% SOURCE_STATEMENT: "Is there an index theoretic interpretation of elliptic genus? When it can, how do we compute it and what is the geometric interpretation of the coefficients? Check that this matches CFT calculations."
% Preparation and final self-audit: CONTRIBUTING.md prompts C and D; source/audit artifacts retained in the local physics_statements_2026_10_11 batch.
% !TeX program = lualatex
\documentclass[11pt]{article}
\usepackage[margin=25mm]{geometry}
\usepackage{fontspec}
\setmainfont{Latin Modern Roman}
\usepackage{amsmath,amssymb,mathtools}
\usepackage{unicode-math}
\setmathfont{Latin Modern Math}
\usepackage{xurl,hyperref}
\hypersetup{colorlinks=true,urlcolor=blue}
\setcounter{secnumdepth}{0}
\setlength{\parindent}{0pt}
\setlength{\parskip}{5pt}
\setlength{\emergencystretch}{3em}
\newcommand{\sref}[2]{\hyperlink{source:#1:#2}{[S#2]}}
\newcommand{\cataloguescope}[1]{\paragraph{Recorded target.}#1}
\newcommand{\R}{\mathbb R}
\newcommand{\C}{\mathbb C}
\newcommand{\N}{\mathbb N}
\newcommand{\Z}{\mathbb Z}
\newcommand{\Q}{\mathbb Q}
\newcommand{\F}{\mathbb F}
\newcommand{\eps}{\varepsilon}
\begin{document}
% UnsolvedMath ID 20002508; source catalogue code AIM-PHYSICS-0016
\setcounter{section}{20002507}
\section{Coefficientwise indices and a level-two Calabi--Yau/CFT test}\label{20002508}
{\small UnsolvedMath ID 20002508 · Mathematical Physics}
\subsection{Definitions and mathematical statement}

\paragraph{Shared notation.}$\mathbb N=\{1,2,\ldots\}$; the standard number fields have their usual meanings. All additional conventions are stated locally.


\paragraph{Definitions and notation.}An index is a signed difference of kernel dimensions of an appropriate Fredholm operator, or its equivariant/graded analogue. An elliptic genus is a generating series of such indices in compact settings. Its Fourier coefficients may count signed or weighted geometric data. CFT denotes conformal field theory, whose state-space trace supplies a separate physical calculation to compare with the geometric series.

\paragraph{Catalogue wording retained for traceability.}
Is there an index theoretic interpretation of elliptic genus? When it can, how do we compute it and what is the geometric interpretation of the coefficients? Check that this matches CFT calculations.
\paragraph{Scope and source qualification.}All three requests are preserved. The source does not prescribe a particular noncompact extension, target Calabi--Yau, grading or CFT; these must be specified in each claimed comparison. The catalogue title/status was paired with a machine-generated research attempt; that attempt is not peer-reviewed evidence of settlement. The statement covers the original source questions rather than only the title’s subcase.
\subsection{Short English statement}
Interpret and compute elliptic genera as indices, explain their coefficients geometrically, and compare with conformal field theory.
\subsection{Sources}
\begin{itemize}
\item[\hypertarget{source:20002508:1}{S1}] ULAM.AI and the UnsolvedMath Contributors, UnsolvedMath, record 20002508 (AIM-PHYSICS-0016). Catalogue curation under CC BY 4.0; definitions and English exposition added locally.
\url{https://huggingface.co/datasets/ulamai/UnsolvedMath}.
\item[\hypertarget{source:20002508:2}{S2}] American Institute of Mathematics, Gromov-Witten invariants and number theory, source item 3.2.
\url{http://aimpl.org/gromwitnumthry/3/}.
\end{itemize}
\label{end:20002508}
\end{document}
